\notechapter{N-20230127}{Cyclotomic Pullbacks: Exact Orders, Galois Conjugacy, and Polynomial Divisibility}

\section{A plausible generalization fails for a precise reason}

Write
\[
  H_m(x)=1+x+x^2+\cdots+x^{m-1}
  =\frac{x^m-1}{x-1}.
\]
The elementary problem
\[
  x^2+x+1\mid x^{2n}+x^n+1
\]
is the case
\[
  H_3(x)\mid H_3(x^n).
\]
For \(m=3\), the answer is simply \(3\nmid n\).  It is tempting to guess that, for general \(m\),
\[
  H_m(x)\mid H_m(x^n)
  \quad\text{whenever}\quad
  m\nmid n.
\]
The guess is false.

Take \(m=6\) and \(n=2\).  Although \(6\nmid2\), the number \(-1\) is a root of
\[
  H_6(x)=1+x+\cdots+x^5,
\]
because
\[
  1-1+1-1+1-1=0.
\]
But
\[
  H_6((-1)^2)=H_6(1)=6,
\]
so \(-1\) is not a root of \(H_6(x^2)\).  Therefore
\[
  H_6(x)\nmid H_6(x^2).
\]

The failed guess looked only at the order \(m\) itself.  The polynomial \(H_m\), however, contains every nontrivial \(m\)-th root of unity, including roots whose exact orders are proper divisors of \(m\).  In the counterexample, the troublesome root has order \(2\), and \(2\mid n\).  The correct question is therefore not merely whether \(m\) divides \(n\), but how the map
\[
  z\longmapsto z^n
\]
changes the exact order of every root of unity involved.

\section{Taking powers changes order by a greatest common divisor}

Let \(\zeta\) be a root of unity of exact order \(d\).  Thus
\[
  \zeta^d=1
\]
and no smaller positive power of \(\zeta\) equals \(1\).

\begin{theorem}[Order after taking a power]
For every positive integer \(n\),
\[
  \operatorname{ord}(\zeta^n)
  =\frac{d}{\gcd(d,n)}.
\]
\end{theorem}

\begin{proof}
Set
\[
  g=\gcd(d,n),
  \qquad
  d=gd_0,
  \qquad
  n=gn_0,
\]
where \(\gcd(d_0,n_0)=1\).  Then
\[
  (\zeta^n)^{d_0}
  =\zeta^{g n_0 d_0}
  =\zeta^{n_0d}=1,
\]
so the order of \(\zeta^n\) divides \(d_0\).

Conversely, if \((\zeta^n)^r=1\), then \(d\mid nr\).  After dividing by \(g\),
\[
  d_0\mid n_0r.
\]
Since \(d_0\) and \(n_0\) are coprime, Euclid's lemma gives \(d_0\mid r\).  Hence no smaller positive exponent than \(d_0\) can kill \(\zeta^n\).
\end{proof}

This formula explains several phenomena at once.  If \(\gcd(d,n)=1\), taking the \(n\)-th power preserves exact order.  If \(d\mid n\), it collapses the root all the way to \(1\).  Between these extremes, the order decreases by exactly the common factor.

For example, let \(\zeta\) have order \(12\).  Then
\[
  \operatorname{ord}(\zeta^8)
  =\frac{12}{\gcd(12,8)}
  =3.
\]
Thus the eighth-power map sends primitive twelfth roots to primitive cubic roots.  It neither preserves all twelve vertices nor collapses them to one point; it folds the twelve-cycle onto a three-cycle.

\section{The correct divisibility criterion}

The roots of \(H_m(x)\) are precisely the nontrivial \(m\)-th roots of unity.  Indeed,
\[
  (x-1)H_m(x)=x^m-1,
\]
and \(x=1\) is the only \(m\)-th root removed by the quotient.

Now let \(\zeta\neq1\) satisfy \(\zeta^m=1\).  Then
\[
  H_m(\zeta^n)=0
\]
exactly when \(\zeta^n\) is still a nontrivial \(m\)-th root of unity.  It is automatically an \(m\)-th root because
\[
  (\zeta^n)^m=(\zeta^m)^n=1.
\]
The only possible failure is \(\zeta^n=1\).

\begin{theorem}[General geometric-sum divisibility]
For positive integers \(m,n\),
\[
  H_m(x)\mid H_m(x^n)
  \quad\Longleftrightarrow\quad
  \gcd(m,n)=1.
\]
\end{theorem}

\begin{proof}
Assume first that \(\gcd(m,n)=1\).  Every root \(\zeta\) of \(H_m\) has exact order \(d>1\) dividing \(m\).  Therefore \(\gcd(d,n)=1\), and the order formula gives
\[
  \operatorname{ord}(\zeta^n)=d>1.
\]
Hence \(H_m(\zeta^n)=0\).  Every root of \(H_m\) is a root of \(H_m(x^n)\), and the roots of \(H_m\) are simple, so divisibility follows.

Conversely, suppose
\[
  g=\gcd(m,n)>1.
\]
Choose a primitive \(g\)-th root of unity \(\eta\).  Since \(g\mid m\), the number \(\eta\neq1\) is a root of \(H_m\).  Since \(g\mid n\),
\[
  \eta^n=1,
\]
and therefore
\[
  H_m(\eta^n)=H_m(1)=m\neq0.
\]
Thus divisibility fails.
\end{proof}

The condition is stronger than \(m\nmid n\) because every nontrivial divisor of \(m\) matters.  Coprimality guarantees that the power map is an automorphism of the entire cyclic group of \(m\)-th roots, not merely that a primitive \(m\)-th root avoids collapsing to \(1\).

\section{The exact common factor is stronger information}

When full divisibility fails, the two polynomials may still share a substantial factor.  Exact orders identify it completely.

Let \(\zeta\) have exact order \(d\), where \(d\mid m\) and \(d>1\).  Then \(\zeta\) is a root of \(H_m(x^n)\) precisely when
\[
  \zeta^n\neq1.
\]
By the order formula, this is equivalent to
\[
  d\nmid n.
\]
Thus the roots shared by \(H_m(x)\) and \(H_m(x^n)\) are exactly the roots whose exact orders divide \(m\) but do not divide \(n\).

To turn this root description into a polynomial formula, group roots by exact order.  Denote by \(\Phi_d(x)\) the polynomial whose roots are the primitive \(d\)-th roots of unity.  Then
\[
  H_m(x)=\prod_{\substack{d\mid m\\d>1}}\Phi_d(x).
\]
Consequently,
\[
  \boxed{
  \gcd\bigl(H_m(x),H_m(x^n)\bigr)
  =
  \prod_{\substack{d\mid m,\ d>1\\d\nmid n}}
  \Phi_d(x).}
\]
The gcd is taken monic in \(\mathbb Q[x]\).

For \(m=6\) and \(n=4\),
\[
  H_6(x)=\Phi_2(x)\Phi_3(x)\Phi_6(x).
\]
The relevant orders are \(2,3,6\).  Since \(2\mid4\) but \(3\nmid4\) and \(6\nmid4\),
\[
  \gcd\bigl(H_6(x),H_6(x^4)\bigr)
  =\Phi_3(x)\Phi_6(x).
\]
Using
\[
  \Phi_3(x)=x^2+x+1,
  \qquad
  \Phi_6(x)=x^2-x+1,
\]
we obtain
\[
  \gcd\bigl(H_6(x),H_6(x^4)\bigr)
  =x^4+x^2+1.
\]
The factor is neither trivial nor all of \(H_6\).  It records exactly which order strata survive the fourth-power map.

\section{Why one primitive root forces an entire rational factor}

A root test over \(\mathbb C\) becomes a divisibility statement over \(\mathbb Q\) only because primitive roots come in algebraically inseparable families.

Let
\[
  \zeta_d=e^{2\pi i/d}
\]
be a primitive \(d\)-th root of unity.  Its powers
\[
  \zeta_d^a,
  \qquad
  a\in(\mathbb Z/d\mathbb Z)^\times,
\]
are exactly the other primitive \(d\)-th roots.  The polynomial
\[
  \Phi_d(x)
  =
  \prod_{a\in(\mathbb Z/d\mathbb Z)^\times}
  (x-\zeta_d^a)
\]
is the \(d\)-th cyclotomic polynomial.

A fundamental theorem says that \(\Phi_d(x)\) lies in \(\mathbb Z[x]\) and is irreducible over \(\mathbb Q\).  Therefore it is the minimal polynomial of \(\zeta_d\) over \(\mathbb Q\).  If a rational polynomial \(P(x)\) satisfies
\[
  P(\zeta_d)=0,
\]
then polynomial division by the minimal polynomial forces
\[
  \Phi_d(x)\mid P(x).
\]
Indeed, write
\[
  P(x)=Q(x)\Phi_d(x)+R(x),
  \qquad
  \deg R<\deg\Phi_d.
\]
Evaluating at \(\zeta_d\) gives \(R(\zeta_d)=0\).  Minimality implies \(R=0\).

The field
\[
  \mathbb Q(\zeta_d)
\]
is the cyclotomic field generated by \(\zeta_d\).  Every automorphism fixing \(\mathbb Q\) is determined by
\[
  \zeta_d\longmapsto\zeta_d^a
\]
for a unit \(a\bmod d\), and every such choice occurs.  Thus
\[
  \operatorname{Gal}(\mathbb Q(\zeta_d)/\mathbb Q)
  \cong(\mathbb Z/d\mathbb Z)^\times.
\]
If \(P\in\mathbb Q[x]\) vanishes at one primitive root, applying these automorphisms shows that it vanishes at every Galois conjugate primitive root.  The minimal-polynomial proof and the Galois proof are two descriptions of the same rigidity.

For \(d=5\), the primitive roots are
\[
  \zeta_5,\zeta_5^2,\zeta_5^3,\zeta_5^4,
\]
and
\[
  \Phi_5(x)=x^4+x^3+x^2+x+1.
\]
A rational polynomial cannot vanish at \(\zeta_5\) alone; it must contain the whole four-root orbit as a factor.

\section{Cyclotomic polynomials separate roots by exact order}

Every \(m\)-th root of unity has a unique exact order \(d\) dividing \(m\).  Partitioning the roots of \(x^m-1\) by exact order gives
\[
  \boxed{x^m-1=\prod_{d\mid m}\Phi_d(x).}
\]
The factors are pairwise coprime because their roots have different exact orders.

Taking degrees yields
\[
  m=\sum_{d\mid m}\deg\Phi_d.
\]
Since the primitive \(d\)-th roots are indexed by the units modulo \(d\),
\[
  \deg\Phi_d=\varphi(d),
\]
so the degree identity becomes the familiar divisor sum
\[
  \sum_{d\mid m}\varphi(d)=m.
\]
The arithmetic identity counts elements of a cyclic group by exact order, while the polynomial identity groups roots by the same invariant.

For \(m=12\), the divisors are
\[
  1,2,3,4,6,12.
\]
The factorization is
\[
\begin{aligned}
x^{12}-1
={}&\Phi_1(x)\Phi_2(x)\Phi_3(x)
\Phi_4(x)\Phi_6(x)\Phi_{12}(x)\\
={}&(x-1)(x+1)(x^2+x+1)(x^2+1)\\
&\cdot(x^2-x+1)(x^4-x^2+1).
\end{aligned}
\]
The degrees
\[
  1+1+2+2+2+4=12
\]
match the twelve roots.  The factor \(\Phi_{12}\) contains exactly the roots of order twelve, not all roots whose orders merely divide twelve.

This exact-order separation is what the original geometric-sum problem needs.  The polynomial \(H_m\) contains every \(\Phi_d\) with \(d\mid m\), \(d>1\), and the power map acts differently on each such stratum.

\section{Substitution pulls cyclotomic factors back along a power map}

To understand \(\Phi_r(x^n)\), ask which roots \(x\) have the property that \(x^n\) has exact order \(r\).  Suppose \(x\) has exact order \(s\).  The order formula says
\[
  \frac{s}{\gcd(s,n)}=r.
\]
Write \(s=rd\).  The condition becomes
\[
  \gcd(rd,n)=d.
\]
This forces \(d\mid n\), and after writing \(n=dq\), it becomes
\[
  \gcd(r,q)=1.
\]
Therefore the possible exact orders of roots of \(\Phi_r(x^n)\) are
\[
  rd
  \quad\text{with}\quad
  d\mid n
  \quad\text{and}\quad
  \gcd\left(r,\frac nd\right)=1.
\]

\begin{theorem}[Cyclotomic pullback formula]
For positive integers \(r,n\),
\[
  \boxed{
  \Phi_r(x^n)
  =
  \prod_{\substack{d\mid n\\
  \gcd(r,n/d)=1}}
  \Phi_{rd}(x).}
\]
\end{theorem}

\begin{proof}
The root analysis above shows that both sides have exactly the same roots, grouped by the exact order \(rd\).  All these roots are simple because \(x^{rn}-1\) has no repeated roots in characteristic zero.  Both sides are monic, so equality follows.
\end{proof}

Take \(r=3\) and \(n=4\).  Every divisor \(d\) of \(4\) satisfies
\[
  \gcd\left(3,\frac4d\right)=1,
\]
so
\[
  \Phi_3(x^4)
  =\Phi_3(x)\Phi_6(x)\Phi_{12}(x).
\]
Indeed,
\[
  x^8+x^4+1
  =(x^2+x+1)(x^2-x+1)(x^4-x^2+1).
\]
The three factors correspond to roots of exact orders \(3,6,12\), all of which become primitive cubic roots after taking fourth powers.

By contrast,
\[
  \Phi_2(x^4)=x^4+1=\Phi_8(x).
\]
Only \(d=4\) contributes because
\[
  \gcd(2,4/d)=1
\]
holds only when \(d=4\).  An eighth root of exact order eight becomes \(-1\) after the fourth power, while roots of orders two or four collapse to \(1\) instead.

\section{Möbius inversion reconstructs the exact-order factor}

The identity
\[
  x^n-1=\prod_{d\mid n}\Phi_d(x)
\]
expresses an ``order divides \(n\)'' object as a product of ``exact order \(d\)'' objects.  Möbius inversion reverses this divisor aggregation.

Recall the Möbius function:
\[
  \mu(n)=
  \begin{cases}
  1,&n=1,\\
  0,&n\text{ is divisible by a square},\\
  (-1)^k,&n\text{ is a product of }k\text{ distinct primes}.
  \end{cases}
\]
Its key property is
\[
  \sum_{d\mid n}\mu(d)
  =
  \begin{cases}
  1,&n=1,\\
  0,&n>1.
  \end{cases}
\]
Multiplicative Möbius inversion gives
\[
  \boxed{
  \Phi_n(x)
  =\prod_{d\mid n}(x^d-1)^{\mu(n/d)}.}
\]
Although negative exponents appear, the total expression is a polynomial because it is obtained by inverting the already established cyclotomic factorization.

The cancellation can be checked factor by factor.  Substitute
\[
  x^d-1=\prod_{e\mid d}\Phi_e(x)
\]
into the right-hand side.  The exponent of a fixed \(\Phi_e(x)\) becomes
\[
  \sum_{\substack{d\mid n\\e\mid d}}
  \mu(n/d).
\]
Writing \(d=eq\), this is
\[
  \sum_{q\mid n/e}\mu\left(\frac{n/e}{q}\right),
\]
which equals \(1\) when \(e=n\) and \(0\) otherwise.  Every unwanted exact-order factor cancels.

For \(n=12\), the nonzero values of \(\mu(12/d)\) occur for
\[
  d=2,4,6,12.
\]
Thus
\[
  \Phi_{12}(x)
  =\frac{(x^{12}-1)(x^2-1)}{(x^6-1)(x^4-1)}.
\]
Using
\[
  x^{12}-1=(x^6-1)(x^6+1)
\]
and
\[
  x^4-1=(x^2-1)(x^2+1),
\]
we obtain
\[
  \Phi_{12}(x)
  =\frac{x^6+1}{x^2+1}
  =x^4-x^2+1.
\]
The formula extracts the primitive twelfth roots from all lower-order roots by inclusion--exclusion on the divisor lattice.

\section{The root plots have an exact formula, not merely a pattern}

Consider
\[
  H_m(x^n)
  =1+x^n+x^{2n}+\cdots+x^{(m-1)n}.
\]
Using the geometric-series identity,
\[
  H_m(x^n)
  =\frac{x^{mn}-1}{x^n-1}.
\]
Hence its roots are the \(mn\)-th roots of unity that are not \(n\)-th roots of unity.

Let
\[
  \xi=e^{2\pi i/(mn)}.
\]
Every root has the form
\[
  \xi^q,
  \qquad
  0\leq q<mn,
\]
subject to
\[
  m\nmid q.
\]
Indeed,
\[
  (\xi^q)^n=e^{2\pi iq/m}=1
\]
exactly when \(m\mid q\).  Therefore
\[
  \boxed{
  \operatorname{Roots}(H_m(x^n))
  =\left\{
  e^{2\pi iq/(mn)}:
  0\leq q<mn,\ m\nmid q
  \right\}.}
\]
There are \(mn-n=n(m-1)\) such roots, matching the degree.

For \(m=6,n=4\), the roots are the twenty fourth roots of unity except
\[
  1,i,-1,-i,
\]
which are precisely the fourth roots of unity.  A plot therefore shows a regular twenty-four-gon with every sixth vertex removed.  The missing pattern is not numerical noise; it is the denominator \(x^4-1\) in
\[
  \frac{x^{24}-1}{x^4-1}.
\]

The exact orders present in the plot can be read from the cyclotomic pullback formula.  Since
\[
  H_m(x^n)
  =\prod_{\substack{r\mid m\\r>1}}\Phi_r(x^n),
\]
each factor \(\Phi_r(x^n)\) splits into the order strata \(rd\) permitted by
\[
  d\mid n,
  \qquad
  \gcd(r,n/d)=1.
\]
Thus the visual rings of angles, the algebraic factorization, and the arithmetic of exact orders are the same data in three forms.

Over \(\mathbb R\), conjugate roots pair into quadratic factors.  If \(\theta\notin\{0,\pi\}\), then
\[
  (x-e^{i\theta})(x-e^{-i\theta})
  =x^2-2\cos\theta\,x+1.
\]
For example,
\[
  \Phi_5(x)
  =\bigl(x^2-2\cos(2\pi/5)x+1\bigr)
   \bigl(x^2-2\cos(4\pi/5)x+1\bigr).
\]
The individual quadratic coefficients lie in \(\mathbb Q(\sqrt5)\), while their product lies in \(\mathbb Z[x]\).  This is the real shadow of the Galois orbit of primitive fifth roots.

\section{Resultants distinguish sharing a root from containing a factor}

The divisibility question asks whether every root of \(H_m\) remains a root after the substitution \(x\mapsto x^n\).  A weaker question asks whether the two polynomials share at least one root.  The resultant detects exactly this weaker condition.

If \(P(x)\) is monic with roots \(\alpha_1,\ldots,\alpha_r\), define
\[
  \operatorname{Res}(P,Q)
  =\prod_{j=1}^{r}Q(\alpha_j).
\]
Equivalently, the resultant is the determinant of the Sylvester matrix of \(P\) and \(Q\).  The product formula shows
\[
  \operatorname{Res}(P,Q)=0
  \quad\Longleftrightarrow\quad
  P,Q\text{ have a common root over }\mathbb C.
\]

Apply this to
\[
  P(x)=H_m(x),
  \qquad
  Q(x)=H_m(x^n).
\]
If \(m\nmid n\), choose a divisor-order stratum that survives.  In fact, a primitive \(m\)-th root \(\zeta\) satisfies
\[
  \zeta^n\neq1
\]
whenever \(m\nmid n\), so
\[
  H_m(\zeta^n)=0.
\]
Thus the resultant is zero.

If \(m\mid n\), then every nontrivial \(m\)-th root \(\zeta\) satisfies
\[
  \zeta^n=1,
\]
so
\[
  H_m(\zeta^n)=H_m(1)=m.
\]
Since \(H_m\) has \(m-1\) roots,
\[
  \operatorname{Res}\bigl(H_m(x),H_m(x^n)\bigr)
  =m^{m-1}.
\]
Therefore
\[
  \boxed{
  \operatorname{Res}\bigl(H_m(x),H_m(x^n)\bigr)
  =
  \begin{cases}
  0,&m\nmid n,\\
  m^{m-1},&m\mid n.
  \end{cases}}
\]

This result sharply separates three levels of interaction:
\[
\begin{array}{c|c}
\text{condition}&\text{polynomial consequence}\\ \hline
m\mid n&\gcd(H_m,H_m(x^n))=1,\\
1<\gcd(m,n)<m&\text{a proper nontrivial common factor may remain},\\
\gcd(m,n)=1&H_m(x)\mid H_m(x^n).
\end{array}
\]
The resultant detects whether the gcd is nontrivial; the exact-order factorization determines the gcd itself; the power-map automorphism criterion determines when the entire divisor survives.

The original factorization problem has therefore grown without changing its central question.  A false congruence guess forced us to track exact orders.  Exact orders forced a decomposition into cyclotomic factors.  Rational divisibility forced us to understand minimal polynomials and Galois conjugacy.  Pulling those factors through \(x\mapsto x^n\) led to a precise factorization theorem, while Möbius inversion reconstructed each exact-order piece.  Finally, the resultant clarified the difference between one surviving orbit and all surviving orbits.  Each new object entered because the previous language could no longer state the next question precisely.
